\chapter{Kết luận}
\label{ch:conclusion}

Chương này tổng kết những đóng góp chính của khóa luận, đối chiếu kết quả đạt
được với các mục tiêu đã đề ra ở Chương~\ref{ch:intro}, đồng thời thảo luận
một số hướng nghiên cứu có thể tiếp tục phát triển trong tương lai.

\section{Tổng kết đóng góp và sản phẩm}
\label{sec:conc-summary}

Mục tiêu của khóa luận là nâng cao khả năng sử dụng mô hình ngôn ngữ nhỏ
trong vai trò web agent thông qua việc xây dựng một quy trình thu thập dữ
liệu huấn luyện tự động từ các website không có nhãn. Để đạt mục tiêu này,
khóa luận đề xuất HSGE, một pipeline khám phá có cấu trúc, trong đó quá
trình tương tác với website được tổ chức dựa trên biểu diễn graph thay vì
chỉ khai thác các task độc lập.

Trong quá trình khám phá, hệ thống phát hiện các skill lặp lại từ trajectory,
kiểm chứng chúng trước khi đưa vào thư viện skill, rồi trừu tượng hóa thành
các operator có precondition và effect tường minh. Trên cơ sở các operator
đã thu nhận, pipeline tiếp tục xây dựng operator graph và ghép chuỗi các
operator tương thích để sinh ra các semantic task nhiều bước. Cách tiếp cận
này cho phép tạo ra những mục tiêu huấn luyện mới dựa trên các năng lực đã
được kiểm chứng thay vì chỉ dựa trên các task gắn với từng trang web riêng
lẻ. Trước khi được sử dụng để fine-tune mô hình, toàn bộ trajectory đều trải
qua các bước đánh giá và tinh lọc tự động nhằm đảm bảo chất lượng dữ liệu,
đồng thời duy trì nguyên tắc \emph{prompt parity} giữa giai đoạn huấn luyện
và suy luận.

Khóa luận tạo ra năm sản phẩm chính như sau.

\begin{enumerate}

\item \textbf{Pipeline HSGE.}
Hệ thống hiện thực đầy đủ quy trình khám phá có cấu trúc, bao gồm khám phá
website theo graph, phát hiện và kiểm chứng skill, induction operator với cơ
chế chuẩn hóa predicate, sinh semantic task từ operator graph và thu thập dữ
liệu huấn luyện. Pipeline có thể vận hành trên cả backend mô phỏng phục vụ
kiểm thử và website thực thông qua trình duyệt.

\item \textbf{Pipeline tinh lọc và tập dữ liệu huấn luyện.}
Khóa luận xây dựng đầy đủ pipeline hậu xử lý dữ liệu gồm judge đánh giá độc
lập các episode, bước viết lại suy luận và trình dựng dataset theo nguyên
tắc prompt parity. Từ đó thu được tập dữ liệu fine-tuning được xây dựng hoàn
toàn tự động trên các website WebArena tự triển khai.

\item \textbf{Mô hình web agent đã fine-tune.}
Qwen3.5-9B được fine-tune bằng LoRA trên tập dữ liệu do pipeline thu thập và
được triển khai bằng vLLM với cùng cấu hình tắt chế độ thinking như trong
quá trình huấn luyện, bảo đảm sự nhất quán giữa huấn luyện và suy luận.

\item \textbf{Đánh giá trên benchmark WebArena.}
Khóa luận xây dựng pipeline đánh giá sử dụng cùng prompt với agent triển khai
thực tế để so sánh mô hình gốc và mô hình sau fine-tune. Trên website
GitLab, mô hình sau fine-tune đạt tỷ lệ thành công
6{,}67\% so với 5{,}56\% của mô hình gốc khi chỉ sử dụng 352 mẫu huấn luyện.
Do mỗi task mới được đánh giá một lần, kết quả này được xem như một tín hiệu
ban đầu về hiệu quả của phương pháp thay vì một kết luận có ý nghĩa thống kê
(Mục~\ref{sec:disc-results}).

\item \textbf{Báo cáo kinh nghiệm triển khai.}
Trong quá trình xây dựng hệ thống, khóa luận tổng hợp các kinh nghiệm và
best practice liên quan đến khám phá website, điều khiển agent, induction
operator, xây dựng dữ liệu và fine-tune mô hình. Các kinh nghiệm này được
trình bày trong Chương~\ref{ch:discussion} cùng hướng dẫn tái lập thực
nghiệm ở Phụ lục~\ref{app:repro}.

\end{enumerate}

Bên cạnh các sản phẩm cụ thể, khóa luận cũng đề xuất một cách tổ chức quá
trình thu thập dữ liệu cho web agent dựa trên biểu diễn graph. Thay vì sinh
trực tiếp các task huấn luyện từ từng trang web, phương pháp trước hết xây
dựng một biểu diễn có cấu trúc về các skill đã được kiểm chứng, sau đó khai
thác biểu diễn này để tạo ra các semantic task nhiều bước. Cách tiếp cận đó
giúp tách biệt quá trình khám phá website với quá trình sinh dữ liệu huấn
luyện, đồng thời mở rộng không gian các task có thể tạo ra từ cùng một tập
trải nghiệm.

\section{Đối chiếu mục tiêu và kết quả}
\label{sec:conc-objectives}

Bảng~\ref{tab:objectives} đối chiếu các mục tiêu đã đề ra trong
Mục~\ref{sec:intro-contrib} với những kết quả đạt được của khóa luận.

\begin{table}[H]
\centering
\caption[Mục tiêu của khóa luận và kết quả đạt được]{Đối chiếu các mục tiêu
đặt ra trong Mục~\ref{sec:intro-contrib} với kết quả đạt được.}
\label{tab:objectives}
\small
\begin{tabular}{@{}L{1.0cm} L{5.6cm} L{6.4cm}@{}}
\toprule
 & \textbf{Mục tiêu} & \textbf{Kết quả} \\
\midrule

O1 &
Thu thập trajectory tự động, bền vững và an toàn từ website không nhãn &
Hoàn thành. Hệ thống hiện thực cơ chế frontier hai tầng kết hợp giới hạn
template, chiến lược duyệt web fail-soft và chế độ an toàn fail-closed,
cho phép pipeline vận hành ổn định trong thời gian dài trên các website
thực (Mục~\ref{sec:method-exploration},
\ref{sec:exp-data}). \\

\addlinespace

O2 &
Sinh semantic task nhiều bước từ graph các skill đã kiểm chứng &
Hoàn thành. Khóa luận xây dựng quy trình phát hiện và kiểm chứng skill,
trừu tượng hóa thành operator có precondition và effect, sau đó sinh
semantic task bằng cách ghép chuỗi các operator tương thích. Việc kiểm
chứng được thực hiện bằng điều kiện tất định hoặc thông qua judge khi cần
thiết (Mục~\ref{sec:method-semantic}). \\

\addlinespace

O3 &
Xây dựng pipeline tinh lọc dữ liệu theo nguyên tắc prompt parity &
Hoàn thành. Pipeline bao gồm bước judge đánh giá độc lập, viết lại chuỗi
suy luận và dựng dataset với cùng cấu trúc prompt sử dụng trong quá trình
suy luận. Các bước thực thi lỗi được loại bỏ và dữ liệu được sắp xếp theo
curriculum trước khi fine-tune
(Mục~\ref{sec:method-data}). \\

\addlinespace

O4 &
Fine-tune và đánh giá một mô hình ngôn ngữ nhỏ trên benchmark chính thức &
Hoàn thành một phần. Mô hình đã được fine-tune và đánh giá trên website
GitLab của WebArena, đạt tỷ lệ thành công tăng từ
5{,}56\% lên 6{,}67\% với 352 mẫu huấn luyện. Tuy nhiên, việc đánh giá trên
các website còn lại và các phép lặp thống kê vẫn đang được tiếp tục
(Mục~\ref{sec:exp-main}--\ref{sec:exp-stats}). \\

\addlinespace

O5 &
Đúc kết các bài học kỹ thuật trong quá trình xây dựng hệ thống &
Hoàn thành. Khóa luận tổng hợp các kinh nghiệm và best practice thu được
trong quá trình phát triển hệ thống, bao gồm khám phá website, điều khiển
agent, induction operator, xây dựng dữ liệu và huấn luyện mô hình
(Chương~\ref{ch:discussion}). \\

\bottomrule
\end{tabular}
\end{table}

\section{Hướng phát triển}
\label{sec:conc-future}

Mặc dù đã đạt được các mục tiêu chính của khóa luận, phương pháp đề xuất vẫn
còn nhiều hướng có thể tiếp tục mở rộng nhằm cải thiện cả quy mô dữ liệu lẫn
hiệu quả của mô hình.

\textbf{Mở rộng quy mô thu thập dữ liệu.}
Kết quả thực nghiệm cho thấy quy mô và độ phủ của dữ liệu vẫn là yếu tố có
ảnh hưởng đáng kể đến chất lượng của mô hình sau fine-tuning. Pipeline HSGE
không phụ thuộc vào một website cụ thể và đã hỗ trợ cơ chế chuyển giao skill
giữa các website có cấu trúc tương đồng. Việc mở rộng quá trình thu thập sang
nhiều website hơn, đồng thời khai thác khả năng tái sử dụng các skill đã học,
có thể tạo ra tập dữ liệu đa dạng và phong phú hơn mà không cần xây dựng lại
toàn bộ quá trình khám phá từ đầu.

\textbf{Khai thác tín hiệu kiểm chứng để huấn luyện.}
Bên cạnh tập dữ liệu dùng cho supervised fine-tuning, pipeline còn tạo ra
nhiều tín hiệu đánh giá như các trajectory đã được judge xác nhận thành công,
các trường hợp thất bại và các episode kết thúc bằng
\texttt{report\_infeasible}. Những tín hiệu này có thể được khai thác trong
các phương pháp tối ưu theo sở thích hoặc học tăng cường nhằm tận dụng cả
những trajectory chưa thành công thay vì chỉ sử dụng các mẫu đúng như trong
SFT.

\textbf{Mở rộng sang các mô hình nhỏ hơn.}
Mặc dù Qwen3.5-9B đã có quy mô nhỏ hơn nhiều mô hình ngôn ngữ lớn đang được
sử dụng cho web agent, chi phí triển khai vẫn còn tương đối cao đối với một
số môi trường thực tế. Một hướng nghiên cứu tiếp theo là chưng cất hành vi
của agent đã fine-tune xuống các mô hình có quy mô từ 1B đến 3B nhằm đánh
giá mức độ duy trì hiệu quả của phương pháp khi dung lượng mô hình tiếp tục
giảm.

\textbf{Khai thác skill graph trong giai đoạn suy luận.}
Trong phạm vi khóa luận, skill graph chỉ được sử dụng để sinh dữ liệu huấn
luyện. Tuy nhiên, cùng một biểu diễn này cũng có thể được khai thác ở thời
điểm suy luận nhằm truy xuất các skill hoặc subgoal liên quan để bổ sung ngữ
cảnh cho mô hình. Điều này mở ra khả năng kết hợp giữa dữ liệu đã học và tri
thức có cấu trúc trong cùng một hệ thống web agent.

\textbf{Mở rộng đánh giá thực nghiệm.}
Một hướng phát triển khác là hoàn thiện đánh giá trên toàn bộ các website của
WebArena với nhiều lần lặp để tăng độ tin cậy thống kê. Ngoài ra, phương pháp
cũng có thể được đánh giá trên các benchmark web agent khác và với các thước
đo chi tiết hơn ngoài tỷ lệ thành công nhị phân nhằm đánh giá đầy đủ hơn khả
năng tổng quát hóa của hệ thống.

\vspace{1em}
\noindent
Khóa luận đã đề xuất HSGE, một pipeline tự động thu thập dữ liệu huấn luyện
cho web agent từ các website không có nhãn thông qua quá trình khám phá có
cấu trúc. Kết quả thực nghiệm ban đầu cho thấy dữ liệu được tạo ra từ
pipeline có thể được sử dụng để cải thiện hiệu quả của một mô hình ngôn ngữ
nhỏ trên benchmark WebArena. Mặc dù vẫn còn nhiều hướng cần tiếp tục nghiên
cứu và đánh giá ở quy mô lớn hơn, phương pháp này cho thấy việc đầu tư vào
chất lượng và cấu trúc của dữ liệu huấn luyện là một hướng tiếp cận khả thi
để nâng cao năng lực của các web agent sử dụng mô hình ngôn ngữ nhỏ.
